\documentclass{beamer}[10]
\usepackage{pgf}
\usepackage[utf8]{inputenc}
\usepackage{beamerthemesplit}
\usepackage{graphics,epsfig, subfigure}
\usepackage{url}
\usepackage{srcltx}
\usepackage{hyperref}

\definecolor{kugreen}{RGB}{0,180,114}%{50,93,61} 
\definecolor{kugreenlys}{RGB}{132,158,139}
\definecolor{kugreenlyslys}{RGB}{173,190,177}
\definecolor{kugreenlyslyslys}{RGB}{0,135,114}%{214,223,216}
\setbeamercovered{transparent}
\mode<presentation>
\usetheme[numbers,totalnumber,compress,sidebarshades]{PaloAlto}
\setbeamertemplate{footline}[frame number]
 
  \usecolortheme[named=kugreen]{structure}
  \useinnertheme{circles}
  \usefonttheme[onlymath]{serif}
  \setbeamercovered{transparent}
  \setbeamertemplate{blocks}[rounded][shadow=true]
\logo{\includegraphics[width=1.5cm]{figs/logo_branca}}
\setbeamercolor{logo}{bg=kugreenlyslyslys}
%\useoutertheme{infolines} 
\title{Turbine Repair}
\author{Gabriel Alcantara}
\institute{LEAD \\ COPPETEC}
\date{05/15/2015}



\begin{document}
\frame{\titlepage \vspace{-0.5cm} 
}

\frame
{
\frametitle{Overview}
\tableofcontents%[pausesection]
}

\section{Cavitation Repair}
\frame{
\frametitle{Cavitation Repair}
\begin{itemize}
  \item The repair of cavitation pitting damage on turbines is an essential part of a
hydro plant maintenance program.
  \item Unrepaired or improperly repaired damages tends to get worse, usually
  at an accelerating rate
\end{itemize}

}

\subsection{Cavitation Repair}

\frame{
\frametitle{Cavitation Repair Steps}
The basic steps of such a repair program are as follows:

\begin{itemize}
  \item Inspection
  \item Identify cause of pitting
  \item Plan best approach to repairs
  \item Perform repairs
\end{itemize}

}


\frame{
\frametitle{Frequency of inspection and repair}

\begin{itemize}
  \item When a new turbine is placed into
operation, the unit should be inspected after about 1500, 4000, and 8000 hours
of operation
\item Turbines . . . . . . . . . . every 24,000 operating hours
or every 4 years
\end{itemize}
}

\frame{
\frametitle{Repair Strategies}

\begin{itemize}
  \item Make all repairs each inspection period
  \item Repair only areas where cavitation damage is 3mm or
deeper
\item Repair areas on stainless steel overlay where pitting is 3 mm
or deeper
\item Allow cavitation to progress to the maximum depth which can be
repaired with two weld passes (about 10mm)
\item Make repairs only when damage becomes so bad that it threatens to
impair the strength of the turbine or when preparation of the damaged
area may result in holes completely through the runner blade
\end{itemize}

}

\frame{
\frametitle{Economic Evalutation}
An economic evaluation of the optimum repair frequency should consider the
following costs:

\begin{itemize}
  \item Unwatering and dewatering the unit
  \item Setting up and removing of the runner maintenance platform
  \item Labor and supervision for repair
  \item Material
  \item Lost revenue due to downtime
  \item \textbf{Possibility to recoat the blades} 
\end{itemize}
}

\subsection{Repair Methods}

\frame{
\frametitle{Repair Methods}
\begin{itemize}
  \item Fill damaged area with weld material
  \item Fill damaged area with non-fused materials
  \item Weld plates over the damaged area
  \item Remove damaged section and replace with new forced plates welded in
place
\end{itemize}
}

\subsubsection{Weld Repair}
\frame{
\frametitle{Weld Repair}

The various steps
for repair by welding are as follows:

\begin{itemize}
  \item Initial dimensional checks
  \item Application of braces of strongbacks Surface preparation
  \item Preheat
  \item Weld application
  \item Grinding to contour
  \item Weld inspection
  \item Removal of braces and strongbacks Final dimensional checks
\end{itemize}
}

\frame{
\frametitle{Surface Preparation}
\begin{itemize}
  \item Grinding
  \item Chipping
  \item Carbon arc gouging
\end{itemize}
}
%“Surface Preparation. Surface preparation should be by grinding, chipping, or
%carbon arc gouging (air arcing) the damaged area to sound metal. Following
%carbon arc gouging, the prepared area should be ground to bright metal.
%Grinding is particularly important on martensitic stainless steel castings to
% assist in controlling the heat-affected zone from the carbon arc gouging, and should
%extend to 2 inches (51 mm) beyond the prepared area. For austenitic stainless
%steel, chipping is not recommended because this material will tend to work
%harden, making preparation more difficult.}

\frame{
\frametitle{Inspection}
The prepared area should be visually inspected for defects.

In high stressed
areas, inspection using:

\begin{itemize}
  \item Magnetic particles (on carbon steel or
martensitic stainless steel)
\item Dye penetrant (on austenitic stainless steel)
procedures should be used)
\end{itemize}
}

\frame{
\frametitle{Preheating}

Preheating is recommended under certain circumstances as part
of good welding practice. 

\begin{itemize}
  \item The purpose of preheating is to prevent hydrogen-induced cracking
\end{itemize}}

\frame{
\frametitle{Weld application}

Weld application should be aimed at achieving good quality weld
with minimum \textbf{residual stressing} and also at keeping \textbf{distortion
to an absolute minimum}. 
}



\frame{
\frametitle{Weld application}

This can be accomplished by:

\begin{itemize}
  \item Less than 3 mm weld material per pass
  \item Limit heat input to the runner
  \item The maximum interpass temperature should
be within 111 C of the minimum preheat temperature
\item Temperature
should be monitored with a contact pyrometer or temperature-indicating
crayons
\item the area about 4 inches (102 mm) from the weld should
not become to warm to touch
\item Apply weld using a cross pattern or a temper bead technique. Cross pattern
welding should be arranged such that the final pass is in the direction of
water flow
\end{itemize}
}
%TODO review this


\frame{
\frametitle{Grinding to Contour}
\begin{itemize}
  \item Following the completion of welding, surfaces should be ground
to their proper profile.
\item Profile should be checked using a flexible spline about 2 feet
(0.6 m) long. For large areas of repair, blade templates should be made.
\end{itemize}
}

\frame{
\frametitle{Tolerances}
The recommended upper limit surface finishes are
as follows:

\begin{center}
  \begin{tabular}{ c | c  }
    \hline
    Head & Surface Finish \\ \hline
    \textless 91 m & 6.3 u m \\ \hline
    91-305 m & 3.2 u m \\ \hline
    305 m & 2.3 u m \\ \hline
    \hline
  \end{tabular}
\end{center}
}

\frame{
\frametitle{Tolerances}

\begin{figure}
\includegraphics[width=0.9\columnwidth]{figs/tolerance}
\end{figure}


}

\frame{
\frametitle{Weld Inspection}


Inspection for small defects or holes in the
stainless steel overlay on a carbon steel runner

}
%TODO Review this

\frame{
\frametitle{Final Dimensional Checks} 

Final dimensional checks should include a
repeat of those dimensions measured prior to the start of repairs.

\begin{itemize}
  \item Comparsion with templates
  \item Vent opening
after the repair should be equal to or greater than that measured prior to the repair
\item A tolerance of -0 to + 3 percent of total vent is recommended
\item On Kaplan runners, the blade angle should also be unchanged. A
tolerance of -0 to + 1/2 degree is recommended
\end{itemize}
}

\subsubsection{Non-Fused Materials}


\frame{
\frametitle{Non-Fused Materials}

Various non-fused materials have been used for
cavitation damage repairs. These include:

\begin{itemize}
  \item Epoxies
  \item Ceramics
  \item Metal spraying
  \item Neoprene
  \item Urethane
\end{itemize} 

The application of these materials is limited by the thickness of the repair.}


\subsubsection{Welding with Solid Plate}
\frame{
\frametitle{Repairs by Welding With Solid Plate}

Cavitation pitting may be repaired by welding solid plate over the
damaged areas 

\begin{itemize}
  \item Repair should be restricted to thick sections
where stresses are low
\item Repair should not be in an area which is subjected
to vibration or dynamic loading
\end{itemize}
}

\frame{
\frametitle{Repairs by Welding With Solid Plate}

\begin{figure}
\includegraphics[width=0.3\columnwidth]{figs/plates}
\end{figure}

\begin{figure}
\includegraphics[width=0.3\columnwidth]{figs/prototyping}
\end{figure}
}

\frame{
\frametitle{Remove Damage Section and Replace with New Plates Welded in Place}

The removal of a severely damaged section of a runner blade and
replacement with a new plate is relatively common in both shop and field repair
of turbine runners.}

\section{Crack Repair}

\frame{
\frametitle{Crack Repair}

Cracks may appear which are caused by localized residual stress or solid objects
striking the blades of the unit.

\begin{itemize}
  \item Records should be kept on the locations and lengths of
these cracks
\item A crack can relieve the residual stress and no further cracking will
develop
\item Holes drilled at the ends of the crack (stop drilling)
\item Welding of runner cracks should not be attempted
\end{itemize}
}
 
 
 %TODO table of repairs
 




\end{document}
